\worksheettitle
  {M01-R. Сколько концов?}
  {\modulemark{M01-R} Коррекционная карточка}

\studentnameblock

\begin{warningbox}[Назовите ошибку]
  Если объект назван только по внешнему виду линии, легко перепутать прямую,
  отрезок и луч. Нужно проверить число концов и направление продолжения.
\end{warningbox}

\begin{practicebox}[Шаг 1. Заполните опору]
  Прямая имеет \answerblank[18mm] концов и продолжается
  \answerblank[35mm]. Отрезок имеет \answerblank[18mm] конца. Луч имеет
  \answerblank[28mm] и продолжается в \answerblank[28mm] сторону.
\end{practicebox}

\begin{practicebox}[Шаг 2. Прочитайте название луча]
  В записи «луч $MN$» начало — точка \answerblank[15mm]. Луч выходит из
  неё и проходит через точку \answerblank[15mm]. Обратный луч называется
  \answerblank[22mm].
\end{practicebox}

\begin{practicebox}[Шаг 3. Исправьте]
  Ученик назвал луч, выходящий из $A$ через $C$, лучом $CA$.

  Причина: \answerblank[85mm]

  Исправление: \answerblank[55mm]
\end{practicebox}

\begin{checkpointbox}[Повторная проверка]
  \begin{enumerate}[label=\textbf{\arabic*.}]
    \item Объект имеет два конца $P,Q$. Это \answerblank[25mm].
    \item Объект начинается в $D$, проходит через $K$ и продолжается
      дальше. Его название: \answerblank[25mm].
    \item Чем отличаются лучи $DK$ и $KD$? \writinglines{1}
  \end{enumerate}
\end{checkpointbox}
